\scanpage{5}
\(\overset{\hspace{0.55em}\unsure{?|}}{\text{\uline{NB}}}\)\quad If $\sigma$ is an $\underset{\hspace{-1.3em}\unsure{\cdot|}}{\text{automorphism}}$ of the field of complex numbers, and if $f$ is ``special'', then $f^\sigma$ is \unsure{also|also} special -- thus $\operatorname{Aut}(\mathbb C)$ acts on the set of special $f$. But it acts through $\operatorname{Aut}(\unsure{\overline{\mathrm Q}|\overline{\mathbb Q}})$, for if $f=P/Q$ special, the $\underset{\hspace{1.5em}\unsure{-|=|}}{\text{coefficients}}$ of $P,Q$ are \uline{algebraic numbers}. It amounts to the same to say that the zeros of $P,Q$ are algebraic numbers, i.e. the divisor of $f$ is

% A pale hook and dark dot above NB are retained as an unsure question-like shape.
% A tiny detached dot below automorphisme is retained; its function is uncertain.
% Faint detached strokes below the right end of coeff. are preserved above;
% the mark's origin and intention remain unresolved (full scan x409--436, y887--902).
% The following drawing preserves the relative placement of four detached
% lower mark-groups. The three uncertain groups have unresolved intention;
% no mathematical relation between them is supplied.
\figure{\begin{tikzpicture}[x=0.011cm,y=0.011cm]
  \node[inner sep=0pt] at (489,161) {$\unsure{\cdot\,\cdot|-}$};
  \node[inner sep=0pt] at (0,0) {$\unsure{-|=}$};
  \node[inner sep=0pt] at (61,5) {$\unsure{\frown|-|}$};
  \node[inner sep=0pt] at (545,0) {$-$};
\end{tikzpicture}}
% archive sheet number: 4
